\section{Both quadratic conductor-eight families in the actual arithmetic window}
\label{sec:quadratic8-actual-window}

\paragraph{Status and dependency boundary.}
This is a conditional-interface proof about the existing arithmetic
window.  It uses the original strict theorems AC-G149 and AC-G150,
with AC-G150-class applied only through the latter, and retains
all the established sectors of AC-G148-window.
The exact author and independent checks are recorded separately.
It proves no estimate for the remaining indices and no mixed
realization.  No additional finite host character at $269$ is used.

Put $p=269$, $K=\mathbb Q(\omega)$ and $t^3=p$ with $t$ real,
and let $\chi$ describe $K(t)/K$.
Let $\phi$ be the primitive odd character modulo sixteen with
$\phi(13)=i$ and $\ker\phi=\{1,7\}$.
Write
\[
 \psi_4=\chi_{-4}|_{G_K},\qquad
 \xi_+=\chi_8|_{G_K},\qquad \xi_-=\chi_{-8}|_{G_K}.
 \tag{Q8W.1}
\]
Use the actual cyclotomic character $c_p$.
Fix a source datum, $0<\varepsilon<1/256$, an allowed page and
its prescribed pair $q,d$.
The fields, cutoff and normalizations remain exactly
\[
 \begin{aligned}
 F_u&=\mathbb Q(\zeta_{2^{h(u)+2}},\zeta_7,
                         \sqrt{-q},\sqrt{-d},\omega),\\
 H_u&=F_u(t),\qquad G_u=\operatorname{Gal}(H_u/K),\quad D_u=|G_u|,\\
 u^{2+8\varepsilon}&\leq2^{h(u)}\leq2u^{2+8\varepsilon},
                   \qquad T=\lfloor u^{399}\rfloor.
 \end{aligned}
 \tag{Q8W.2}
\]
Use the original $O_\ell,E_\ell,b_\ell$, strict conditions,
derived indices and unrestricted coefficient-prime range of
(FC.1)--(FC.2), (EU.2) and (TC.1).

For $u\geq2$, retain the four disjoint sets of AC-G148:
\[
 \begin{aligned}
 Z_{16}^+&=\{(p,\chi^{-j}\phi^{-\sigma},e):
      j=1,2,\ \sigma=\pm1,\\
 &\hspace{2em}1\leq e<T,\ e\equiv-1\pmod{67}\},\\
 Z_4&=\{(p,\chi^{-j}\psi_4,e):
      j=1,2,\ 1\leq e<T,\ e\equiv-1\pmod{134}\},\\
 Z_{\rm cub}&=\{(p,\chi^j,e):
      j=0,1,2,\ 1\leq e<T,\ e\equiv0,267\pmod{268}\},\\
 Z_3&=\{(3,1,e):1\leq e<T\}.
 \end{aligned}
 \tag{Q8W.3}
\]
Adjoin the new set
\[
 \begin{aligned}
 Z_8&=\{(p,\chi^{-j}\xi_\nu,e):j=1,2,\ \nu\in\{+,-\},\\
 &\hspace{2em}1\leq e<T,\ e\equiv-1\pmod{134}\},\\
 Z^{+,8}&=Z_{16}^+\sqcup Z_4\sqcup Z_{\rm cub}\sqcup Z_3\sqcup Z_8.
 \end{aligned}
 \tag{Q8W.4}
\]
Let $I_{399}^{\rm est8}$ be the sum (TC.1) with precisely
$Z^{+,8}$ omitted, and let $I_{399}^{\rm est+}$ denote the
AC-G148 remaining sum.  Every unlisted prime, finite character
and positive Tate weight is retained.

\begin{theorem}[AC-G151-window: the quadratic-eight actual-label reduction]
\label{AC-G151-window}
All labels in (Q8W.3)--(Q8W.4) occur on each prescribed host,
and their strict groups on the original coefficients (FC.2) vanish.
The exact disjoint counts, for the positive integer $T$, are
\[
 \begin{aligned}
 |Z_{16}^+|&=4\lfloor T/67\rfloor,&
 |Z_4|&=2\lfloor T/134\rfloor,\\
 |Z_8|&=4\lfloor T/134\rfloor,&
 |Z_3|&=T-1,\\
 |Z_{\rm cub}|&=3\bigl(\lfloor(T-1)/268\rfloor
                                  +\lfloor T/268\rfloor\bigr).
 \end{aligned}
 \tag{Q8W.5}
\]
Thus exactly $4\lfloor T/134\rfloor$ further labels are omitted
relative to AC-G148.
If $C_u$ is the sum of the actual nonnegative normalized
positive-index costs over the whole union $Z^{+,8}$, then
\[
 \begin{aligned}
 I_{399}&=I_{399}^{\rm est8}+C_u,\qquad
       I_{399}^{\rm est8}\leq I_{399}^{\rm est+}\leq I_{399},\\
 0\leq C_u&\leq8D_u^2T^2\log(8TD_u)
                  +\tfrac12D_u\mathcal R_uT^2+2D_ur_uc_0T\\
 &=O_\varepsilon\bigl(u^{802+16\varepsilon}(\log u)^2\bigr).
 \end{aligned}
 \tag{Q8W.6}
\]
The equality is also an identity of nonnegative extended sums;
no subtraction of infinite quantities is used.
Only one full (FL.4) comparison cost is charged for the entire union.

The common-window premise (TC.3), or its AC-G148 form, can therefore
use $I_{399}^{\rm est8}$ with all simultaneous field, coefficient,
character and Tate-weight quantifiers unchanged.
Either old premise implies the new one with its original parameters.
The new premise implies the full old one on the same $W,J$, datum,
page and pair, after increasing only the starting threshold and
halving the positive relative slack.
The mixed theta/Selmer and analytic input $\mathcal R_{\rm cub}$,
the actual source $R.t$, and every other AC-K35 hypothesis remain.
Neither that input nor the remaining numerical bound is proved here.
\end{theorem}

\begin{proof}
\proofstep{1. Four distinct new labels on the same fields.}
The page lower bound and $u\geq2$ imply $h(u)\geq2$.
Hence $F_u$ contains $\mathbb Q(\zeta_{16})$, in particular
$\mathbb Q(\zeta_8)$, and $H_u$ contains $K(t)$.
The characters $\xi_+$ and $\xi_-$ factor through $G_u$.
The fields $K$ and $\mathbb Q(\zeta_8)$ intersect in $\mathbb Q$:
the three quadratic subfields of the latter are
$\mathbb Q(i),\mathbb Q(\sqrt2),\mathbb Q(\sqrt{-2})$,
none of which is $K$.
Thus $1,\psi_4,\xi_+,\xi_-$ remain four distinct characters
on $G_K$ and on the corresponding quotient of $G_u$.

Since $F_u/\mathbb Q$ is abelian, it cannot contain the
nonabelian $S_3$ splitting field $K(t)$.
The cyclic cubic extension $K(t)/K$ has prime degree, so
$F_u\cap K(t)=K$.
Its cubic subgroup distinguishes the two values of $j$.
The two new quadratic factors are distinct from the factor $\psi_4$
of $Z_4$ and the trivial factor of $Z_{\rm cub}$.
The conductor-sixteen factors have order four, so cannot equal
these quadratic factors; $Z_3$ has a different coefficient prime.
This proves all required disjointness and occurrence assertions.

The original coefficient ring $O_p$ contains the cubic values,
since $\exp(G_u)=3\cdot2^{h(u)}$; the new quadratic values are
already rational.  The original $b_p$ is not changed.
For every new label the coefficient is literally
\[
 \begin{aligned}
 A_{j,\nu,e}&=(O_p\otimes_{\mathbb Z_p}\mathbb Q_p/\mathbb Z_p)
                (\chi^{-j}\xi_\nu)(-e)\\
       &=(E_p/O_p)(\chi^{-j}\xi_\nu c_p^{-e}).
 \end{aligned}
 \tag{Q8W.7}
\]
For $\nu=+$ apply AC-G149; for $\nu=-$ apply AC-G150.
In both cases set the additional exponent $m=0$ and continuous
character $\kappa=1$.
Both theorems allow exactly this original $O_p$, with no absolute
ramification restriction, and both residues $133,267\bmod268$.
Their auxiliary quadratic $\tau_p$ belongs only to residual
Kummer descent at residue $133$.
It does not appear in (Q8W.7) or as a new finite host label.
All old strict assertions on (Q8W.3) are retained from AC-G148.

\proofstep{2. Preserve full local conditions and count the union.}
The two new strict theorems use zero in full
$H^1(G_{K_v},A_{j,\nu,e})$ for every $v\mid p$, and zero
in $H^1(I_v,A_{j,\nu,e})$ at every other finite place.
These are the local conditions in (FC.2).
For clarity about its finite-set presentation, (Q8W.7) is
unramified away from $2,3,p$.
At an added away-$p$ place of $\Sigma_p$, inertia acts trivially.
A strict inertia restriction is a coboundary, hence is zero.
The cocycle identity then makes it zero on every conjugate inertia
and on the closed normal subgroup they generate.
The action and cocycle descend to the smaller ramification set.
Conversely inflation is injective and preserves the remaining
local conditions.  This identifies the strict kernels without
losing a full local condition at $p$.
The corresponding original-coefficient comparisons for (Q8W.3),
including the coefficient-three comparison without averaging by
three, are retained from AC-G148 and AC-G146.

The two residues for each new finite label combine exactly to
$e=134n-1$, $n\geq1$.
For positive integer $T$, the inequalities $1\leq134n-1<T$
are equivalent to $1\leq n\leq\lfloor T/134\rfloor$.
There are four distinct new finite labels, proving the $Z_8$ count.
The other four counts are unchanged AC-G148 identities.
Thus precisely these strict terms in (FC.15) vanish.
No conclusion that an individual derived index is zero is used.

\proofstep{3. One local comparison bound for all omitted terms.}
For a removed label $(\ell,\eta,e)$ take the original dual complex
\[
 C_{\ell,\eta,e}=R\Gamma(G_{K,\Sigma_\ell},
                               O_\ell(\eta^{-1})(e+1)).
 \tag{Q8W.8}
\]
The codifferent convention is exactly (FL.2).
The determinant inequality (EU.2) and the proved strict zero give
\[
 \begin{aligned}
 \frac{\log\#k_{O_\ell}}{b_\ell}\max(0,i_{\ell,\eta,e})
 &\leq\frac{\log\#H^2(C_{\ell,\eta,e})}{b_\ell}\\
 &=\frac{\left|\log\#H^2(C_{\ell,\eta,e})
            -\log\#\mathcal S_{\ell,\eta}(e)\right|}{b_\ell}.
 \end{aligned}
 \tag{Q8W.9}
\]
This uses $\log\#\mathcal S=0$ and keeps the positive index cost;
it does not identify strict cohomology with $H^2$.
At a fixed weight, the sum of (Q8W.9) over the entire removed set,
including the coefficient-three labels, is one subsum of the
nonnegative all-prime comparison (FL.4).
It is therefore at most
\[
            16eD_u^2\log(8eD_u)+D_ue\mathcal R_u+2D_ur_uc_0.
 \tag{Q8W.10}
\]
Summing this once over $1\leq e<T$, using
$\sum e\leq T^2/2$, $\log(8eD_u)\leq\log(8TD_u)$ and
$T-1\leq T$, gives (Q8W.6).
No separate full error is added for the new families or any old one.
Nonnegativity gives both additivity and monotonicity, even if
one of the unbounded remaining sums is infinite.

The retained host bounds are $D_u\leq288u^{2+8\varepsilon}$,
$r_u\leq12$, and
$\mathcal R_u=O_\varepsilon(u^{\varepsilon/5}(\log u)^2)$,
uniformly on the prescribed page and pair.
With $T\leq u^{399}$ the first error exponent is
$802+16\varepsilon$, the second $800+41\varepsilon/5$,
and the third is smaller.
For every fixed positive $\varepsilon$ the total error is
$o(\exp(u^\varepsilon/2))$ uniformly in the retained parameters.

\proofstep{4. Keep the same common window and realization.}
Fix the datum, page and $0<\varepsilon<1/256$.
Suppose the new premise supplies $J\geq1$, $X_0$ and
$s\in(0,1)$, independent of $X,q,d,u,\ell,\eta,e$, such that
for every $X\geq X_0$ and prescribed
$\operatorname{WinPair}(\varepsilon,X^{1/J},q,d)$ it supplies
one $W$ satisfying $X^{1/J}\leq W$, $W^2\leq X$ and
\[
 I_{399}^{\rm est8}(H_u,u)\leq(1-s)\exp(u^\varepsilon/2)
                       \qquad(W\leq u\leq W^2).
 \tag{Q8W.11}
\]
Choose $U_0\geq2$ so that the bound in (Q8W.6) is at most
$(s/2)\exp(u^\varepsilon/2)$ for all $u\geq U_0$.
It depends only on the same datum, page, $\varepsilon$ and slack,
not on a new pair, coefficient, finite label, Tate weight or window.
Replace $X_0$ by $\max(X_0,U_0^J)$.
The very same $W$ now satisfies
\[
 I_{399}(H_u,u)\leq(1-s/2)\exp(u^\varepsilon/2)
                           \qquad(W\leq u\leq W^2).
 \tag{Q8W.12}
\]
For a fixed finite simultaneous family, take the maximum of its
uniform error constants, retaining this one $W$.
No intersection of separately chosen windows is invoked.
Conversely the old full or AC-G148 premise implies (Q8W.11)
by nonnegativity with unchanged parameters.
The corresponding uniform little-$o$ formulations agree as well.
The actual source host, rank $R.t$, family construction and every
mixed theta/Selmer and analytic clause of $\mathcal R_{\rm cub}$
remain unchanged.  AC-K35 applies only with those retained inputs.
\end{proof}

\paragraph{Exact old-gate reduction and geometric scope.}
The strongest existing AC-G148 form of the old EB-Q02 premise
still includes the four labels in $Z_8$.
For $T\geq134$, their displayed count is positive, so the new
unknown sum omits a proper additional subset of actual labels.
This is an inclusion statement about terms requiring a numerical
bound, not a claim that the two numerical sums are strictly unequal.
All prior omissions are retained without being counted as new results.
The remaining all-prime index bound, mixed realization, analytic
comparison and H04.4 remain open.
The coefficients are the original rank-one Artin--Tate systems
on the same arithmetic curve and finite cover.
SGA~1, Expos\'e~V, \S\S4 and 7 supplies only the previously
retained finite-cover dictionary.
No new motive, primary-source proof, reciprocity law or Lean
execution is asserted.
